\clearpage
\begin{proofbox}
\textbf{\textcolor{CProof}{方法一：最容易懂——数字接龙（理解为什么）}}
\begin{center}\begin{tikzpicture}[font=\small,>=stealth,line join=round]
\draw[gray!40,rounded corners=3pt] (0,0) rectangle (14.2,8.5);\node at (7.1,7.9) {把输出当作数：双曲线上的两点交换横纵坐标};
\node at (2.4,6.9) {$v=1/u$};\begin{scope}[shift={(2.4,4.4)},x=0.425cm,y=0.45cm]\draw[->,gray] (-4,0)--(4,0) node[right]{$u$};\draw[->,gray] (0,-4)--(0,4) node[above]{$v$};\draw[gray,thick,dashed,domain=-3.6:3.6,samples=160] plot(\x,{\x});\draw[blue,thick,solid,domain=-3.6:-0.28,samples=160] plot(\x,{1/\x});\draw[blue,thick,solid,domain=0.28:3.6,samples=160] plot(\x,{1/\x});\fill[orange!85!black] (2,0.5) circle(2pt);\fill[orange!85!black] (0.5,2) circle(2pt);\end{scope}\node[align=center,text width=4cm] at (2.4,1.7) {$2\to\frac12\to2$};\node at (7.050000000000001,6.9) {$v=-1/u$};\begin{scope}[shift={(7.050000000000001,4.4)},x=0.425cm,y=0.45cm]\draw[->,gray] (-4,0)--(4,0) node[right]{$u$};\draw[->,gray] (0,-4)--(0,4) node[above]{$v$};\draw[gray,thick,dashed,domain=-3.6:3.6,samples=160] plot(\x,{\x});\draw[violet,thick,solid,domain=-3.6:-0.28,samples=160] plot(\x,{-1/\x});\draw[violet,thick,solid,domain=0.28:3.6,samples=160] plot(\x,{-1/\x});\fill[orange!85!black] (2,-0.5) circle(2pt);\fill[orange!85!black] (-0.5,2) circle(2pt);\end{scope}\node[align=center,text width=4cm] at (7.050000000000001,1.7) {$2\to-\frac12\to2$};\node at (11.7,6.9) {$v=-6/u$};\begin{scope}[shift={(11.7,4.4)},x=0.425cm,y=0.45cm]\draw[->,gray] (-4,0)--(4,0) node[right]{$u$};\draw[->,gray] (0,-4)--(0,4) node[above]{$v$};\draw[gray,thick,dashed,domain=-3.6:3.6,samples=160] plot(\x,{\x});\draw[teal,thick,solid,domain=-3.6:-1.67,samples=160] plot(\x,{-6/\x});\draw[teal,thick,solid,domain=1.67:3.6,samples=160] plot(\x,{-6/\x});\fill[orange!85!black] (3,-2) circle(2pt);\fill[orange!85!black] (-2,3) circle(2pt);\end{scope}\node[align=center,text width=4cm] at (11.7,1.7) {$3\to-2\to3$};
\node[text width=12.6cm,align=center] at (7.1,.65) {横轴u为变换前的数，纵轴v为变换后的数；虚线v=u，两点关于它对称。};
\end{tikzpicture}\end{center}



先把$f(x)$暂时看作一个不知道的、但不是$0$的数，记作$y$。已知每向右走$a$，值就按同一个规则改变：
\[
\begin{array}{c|ccc}
 \text{所在位置}&x&x+a&x+2a\\ \hline
 k=1&y&\frac1y&\frac1{1/y}=y\\[0.4em]
 k=-1&y&-\frac1y&-\frac1{-1/y}=y\\[0.4em]
 \text{一般}\ k\ne0&y&\frac{k}{y}&\frac{k}{k/y}=y
\end{array}
\]
两次之后，无论最开始的$y$具体是多少，都能回到$y$。所以相距$2a$的两个位置，函数值总相等。若$a<0$，改用正的长度$2|a|$描述周期即可。

\textbf{这一步解决的疑问：}为什么不是“一走$a$就重复”？因为走第一步通常变成另外一个值，走完第二步才回来。

\medskip

\end{proofbox}
\clearpage
\begin{proofbox}
\textbf{\textcolor{CProof}{方法二：最标准——换一次$x$，连代两次（考试首选）}}
\begin{center}\begin{tikzpicture}[font=\small,>=stealth,line join=round]
\draw[gray!40,rounded corners=3pt] (0,0) rectangle (14.2,8.3);\node at (7.1,7.700000000000001) {一次互为倒数，两次回到同一高度};
\begin{scope}[shift={(1.7923076923076926,2.0951612903225807)},x=1.153846153846154cm,y=1.4838709677419355cm]\draw[->,gray] (-0.6,0)--(4.6,0) node[right]{$x$};\draw[->,gray] (0,-0.3)--(0,2.8) node[above]{$y$};\draw[blue,thick,solid,domain=-0.3:4.3,samples=180] plot(\x,{2^(sin(deg(pi*\x)))});\fill[orange!85!black] (0.5,2) circle(2.2pt);\fill[orange!85!black] (1.5,0.5) circle(2.2pt);\fill[orange!85!black] (2.5,2) circle(2.2pt);\draw[gray,dashed] (-.5,1)--(4.4,1);\node[left] at (0,1){$1$};\foreach \v/\lab in {.5/\frac12,1.5/\frac32,2.5/\frac52}{\node[below=5pt] at (\v,0){$\lab$};}\end{scope}\node[anchor=west,align=left,text width=5.0cm] at (8.4,4.0) {$f=2^{\sin(\pi x)},\ a=1$\\[9pt]$f(\frac12)=2$\\[9pt]$f(\frac32)=\frac12$\\[9pt]$f(\frac52)=2$\\[9pt]正周期可取2};
\node[text width=12.6cm,align=center] at (7.1,.65) {三点高度是2、二分之一、2；倒数不同于把高度变为相反数。};
\end{tikzpicture}\end{center}



先证更一般的情况。设
\[
 f(x+a)=\frac{k}{f(x)},\qquad k\ne0,\ f(x)\ne0.
\]
\begin{description}[style=nextline,leftmargin=2.8em,labelsep=0.5em,itemsep=0.55em]
\item[\textbf{第1步：把$x$换成$x+a$。}] 既然原关系对所有$x$成立，它对$x+a$当然也成立，所以
\[
 f((x+a)+a)=\frac{k}{f(x+a)},\quad
 \text{即}\quad f(x+2a)=\frac{k}{f(x+a)}.
\]
\item[\textbf{第2步：把已知的$f(x+a)$填回去。}]
\[
 f(x+2a)=\frac{k}{\dfrac{k}{f(x)}}
 =k\cdot\frac{f(x)}{k}=f(x).
\]
这里能约去$k$，正是因为题目要求$k\ne0$。
\item[\textbf{第3步：把等式翻译成周期。}]
无论$x$取什么数，都有$f(x+2a)=f(x)$，因此$2|a|$是函数的一个正周期。
\end{description}

如果题目只给\(f(x+a)=-1/f(x)\)，也能逐行写：
\[
\begin{aligned}
 f(x+2a)&=-\frac1{f(x+a)}\\
 &=-\frac1{-1/f(x)}\\
 &=f(x).
\end{aligned}
\]

\medskip

\end{proofbox}
\clearpage
\begin{proofbox}
\textbf{\textcolor{CProof}{方法三：步骤最少——两次乘积相等（速答首选）}}
\begin{center}\begin{tikzpicture}[font=\small,>=stealth,line join=round]
\draw[gray!40,rounded corners=3pt] (0,0) rectangle (14.2,8.3);\node at (7.1,7.700000000000001) {证明三：面积相同、底边相同，所以高相同};
\draw[step=.5,gray!12] (3.4,2.1) grid (4.4,6.1);\draw[step=.5,gray!12] (9.2,2.1) grid (10.2,6.1);\filldraw[fill=blue!12,draw=blue,thick] (3.4,2.1) rectangle (4.4,6.1);\filldraw[fill=orange!12,draw=orange!85!black,thick] (9.2,2.1) rectangle (10.2,6.1);\node[above] at (3.9,6.2){面积1};\node[above] at (9.7,6.2){面积1};\node[left] at (3.3,4.1){$f(x)=2$};\node[right] at (10.3,4.1){$f(x+2)=2$};\node[below=6pt] at (3.9,2.1){$f(x+1)=\frac12$};\node[below=6pt] at (9.7,2.1){$f(x+1)=\frac12$};
\node[text width=12.6cm,align=center] at (7.1,.65) {图取 $a=k=1$；两矩形同尺度，共同底边 $f(x+1)\ne0$ 才能消去。};
\end{tikzpicture}\end{center}



由$f(x+a)=k/f(x)$，两边乘以$f(x)$得到
\[
\underbrace{f(x)f(x+a)}_{\text{第1个乘积}}=k.
\tag{1}
\]
把$x$换成$x+a$，又得到
\[
\underbrace{f(x+a)f(x+2a)}_{\text{第2个乘积}}=k.
\tag{2}
\]
因此
\[
 f(x)f(x+a)=f(x+a)f(x+2a).
\]
因为$f(x+a)\ne0$，两边同时除以它：
\[
\boxed{f(x)=f(x+2a)}.
\]
\textbf{这一招的妙处：}两个乘积都等于$k$，还都有同一个因子$f(x+a)$，一消就结束；几乎不需要处理“分式套分式”。

\medskip
\end{proofbox}\clearpage\begin{proofbox}\textbf{\textcolor{CProof}{补充：为什么负的$a$也能写出正周期？}}
\begin{center}\begin{tikzpicture}[font=\small,>=stealth,line join=round]
\draw[gray!40,rounded corners=3pt] (0,0) rectangle (14.2,8.3);\node at (7.1,7.700000000000001) {负的步长向左取倒数，正周期仍取4};
\begin{scope}[shift={(4.1,2.2075757575757575)},x=0.625cm,y=1.3939393939393938cm]\draw[->,gray] (-4.8,0)--(4.8,0) node[right]{$x$};\draw[->,gray] (0,-0.4)--(0,2.9) node[above]{$y$};\draw[blue,thick,solid,domain=-4.4:4.4,samples=180] plot(\x,{2^(cos(deg(pi*\x/2)))});\fill[orange!85!black] (-4,2) circle(2.2pt);\fill[orange!85!black] (-2,0.5) circle(2.2pt);\fill[orange!85!black] (0,2) circle(2.2pt);\fill[orange!85!black] (4,2) circle(2.2pt);\foreach \v in {-4,-2,0,4}{\ifnum\v=0\node[below left=5pt] at (\v,0){$\v$};\else\node[below=5pt] at (\v,0){$\v$};\fi}\end{scope}\node[anchor=west,align=left,text width=5.0cm] at (8.4,4.0) {$a=-2,\ k=1$\\[9pt]$f=2^{\cos(\pi x/2)}$\\[9pt]$f(-4)=f(0)=f(4)=2$\\[9pt]$f(-2)=\frac12$\\[9pt]$T=4=2|a|$};
\node[text width=12.6cm,align=center] at (7.1,.65) {向左移动4能重复，也就能向右移动4重复；正周期不能写成负4。};
\end{tikzpicture}\end{center}



例如$a=-2$，上面证明的是$f(x-4)=f(x)$。让等式中的$x$换成$x+4$，得到$f(x)=f(x+4)$。所以$4=2|a|$就是正周期。写“周期$T=-4$”不符合“正周期”的通常说法。

\medskip
\textbf{\textcolor{CProof}{最后要分清证明到哪一步}}

三种方法都证明了\(f(x+2a)=f(x)\)，所以证明的是\textbf{$2|a|$是一个周期}，\textbf{没有证明它一定最小}。要求最小正周期时，还必须利用其他信息排除更小的正周期。
\end{proofbox}
